\documentclass[11pt,a4paper]{article}

\usepackage[T1]{fontenc}
\usepackage[utf8]{inputenc}
\usepackage{lmodern}
\usepackage{amsmath,amssymb,mathtools}
\usepackage{booktabs,longtable,array}
\usepackage{geometry}
\usepackage{xcolor}
\usepackage{hyperref}
\usepackage{microtype}

\geometry{margin=2.5cm}
\hypersetup{colorlinks=true,linkcolor=blue!55!black,citecolor=blue!55!black,
  urlcolor=blue!55!black}
\setlength{\parindent}{0pt}
\setlength{\parskip}{0.55em}

\newcommand{\Xist}{\mathrm{Xist}}
\newcommand{\txa}{\mathrm{tXA}}
\newcommand{\cxr}{\mathrm{cXR}}
\newcommand{\Pre}{\mathrm{Pre}}
\newcommand{\Post}{\mathrm{Post}}
\newcommand{\ind}{\mathbf{1}}

\title{Translation of the M\"utzel cXR--tXA Model\\
from Ordinary Differential Equations to a Stochastic Petri Net}
\author{Hadi Fouani\\MPRI M2 -- INRIA MUSCA}
\date{\today}

\begin{document}
\maketitle

\begin{abstract}
This note gives a complete translation of the two-allele cXR--tXA model of
M\"utzel et al.\ from its ordinary differential equation (ODE) formulation to
the stochastic Petri net (SPN) implemented in this project.  The notation
$x_i$, $\txa_i$, $\cxr_i$, and $p_3,\ldots,p_{23}$ follows the paper.  We
explain how each positive or negative ODE term becomes a transition, how
activators and inhibitors determine the choice of input places, how
complementary places impose finite molecular capacities, and how transition
propensities preserve the kinetic functions of the ODE.  The current
implementation splits the production of each $x_i$ into two transitions, one
attributed to each trans-acting activator, without changing their total rate.
\end{abstract}

\tableofcontents

\section{Biological and mathematical model}

Random X-chromosome inactivation requires a cell with two initially active X
chromosomes to reach a state in which one allele expresses high Xist and the
other remains low.  M\"utzel et al.\ classify X-linked regulators by whether
they activate or repress Xist and whether they act in cis or in trans
\cite{mutzel2019}.  The minimal model considered here contains, for each
chromosome $i\in\{1,2\}$,
\begin{itemize}
  \item $x_i$: the level of Xist produced from chromosome $i$;
  \item $\cxr_i$: a cis-acting Xist repressor that affects only allele $i$;
  \item $\txa_i$: the contribution of chromosome $i$ to a trans-acting Xist
        activator shared by both alleles.
\end{itemize}

The regulatory signs are
\[
  \txa_1+\txa_2 \longrightarrow x_i,
  \qquad
  \cxr_i \dashv x_i,
  \qquad
  x_i \dashv \txa_i,
  \qquad
  x_i \dashv \cxr_i.
\]
Thus the tXA pool activates both Xist alleles, whereas cXR represses Xist in
cis.  Once Xist rises on one chromosome, it represses the production of both
local regulators.  The loss of tXA contributes a trans feedback between the
two chromosomes; the loss of cXR creates a positive cis feedback on the same
chromosome.

\section{ODE formulation in the notation of M\"utzel et al.}

For $i\in\{1,2\}$, the stochastic-scale equations reported for the cXR--tXA
model are \cite{mutzel2019supp}
\begin{align}
  \frac{dx_i}{dt}
    &=p_{21}f(\cxr_i)f(\txa)-0.1733x_i, \label{eq:xist}\
  \frac{d\txa_i}{dt}
    &=p_{23}\left(1-
      \frac{x_i^{p_3}}{x_i^{p_3}+(p_{21}p_4)^{p_3}}\right)-\txa_i,
      \label{eq:txa}\
  \frac{d\cxr_i}{dt}
    &=p_{22}\left(1-
      \frac{x_i^{p_5}}{x_i^{p_5}+(p_{21}p_6)^{p_5}}\right)-\cxr_i.
      \label{eq:cxr}
\end{align}
The regulatory functions in the Xist equation are
\begin{align}
 f(\cxr_i)
   &=1-\frac{\cxr_i^{p_{13}}}
      {\cxr_i^{p_{13}}+(p_{22}p_{14})^{p_{13}}}, \label{eq:fcxr}\\
 f(\txa)
   &=\frac{\left[\tfrac12(\txa_1+\txa_2)\right]^{p_{11}}}
      {\left[\tfrac12(\txa_1+\txa_2)\right]^{p_{11}}
       +(p_{23}p_{12})^{p_{11}}}. \label{eq:ftxa}
\end{align}

Every right-hand side has the form
\[
  \frac{dy}{dt}=P_y(\boldsymbol{x})-D_y(\boldsymbol{x}),
\]
where $P_y\geq0$ is production and $D_y\geq0$ is degradation.  This
nonnegative decomposition is the key to the SPN translation: production and
degradation are represented by distinct stochastic events.

The parameters used by the implementation are
\begin{center}
\begin{tabular}{@{}lll@{}}
\toprule
Parameter & Value & Role \\
\midrule
$p_{21}$ & $347$ & Xist molecular scale / maximal production rate \\
$p_{22}$ & $76$  & cXR molecular scale / maximal production rate \\
$p_{23}$ & $79$  & tXA molecular scale / maximal production rate \\
$p_3$    & $3.4$ & Hill exponent for $x_i\dashv\txa_i$ \\
$p_4$    & $0.017$ & normalized threshold for $x_i\dashv\txa_i$ \\
$p_5$    & $2.2$ & Hill exponent for $x_i\dashv\cxr_i$ \\
$p_6$    & $0.019$ & normalized threshold for $x_i\dashv\cxr_i$ \\
$p_{11}$ & $2.7$ & Hill exponent for $\txa\rightarrow x_i$ \\
$p_{12}$ & $1.03$ & normalized threshold for $\txa\rightarrow x_i$ \\
$p_{13}$ & $2.6$ & Hill exponent for $\cxr_i\dashv x_i$ \\
$p_{14}$ & $0.20$ & normalized threshold for $\cxr_i\dashv x_i$ \\
$\delta_x$ & $0.1733\,\mathrm{h}^{-1}$ & Xist degradation constant \\
\bottomrule
\end{tabular}
\end{center}

\section{Discrete places and complementary places}

Let $M(p)$ denote the marking of place $p$.  Each biological variable is
represented by a positive place and a complementary place:
\[
\begin{array}{lll}
 x_i, & \overline{x}_i, & M(x_i)+M(\overline{x}_i)=p_{21},\\
 \txa_i, & \overline{\txa}_i, &
 M(\txa_i)+M(\overline{\txa}_i)=p_{23},\\
 \cxr_i, & \overline{\cxr}_i, &
 M(\cxr_i)+M(\overline{\cxr}_i)=p_{22}.
\end{array}
\]
The source code writes the complementary names as \texttt{!x1},
\texttt{!txa1}, and \texttt{!cxr1}, while this document uses an overline.
One production firing transfers one token from $\overline y$ to $y$; one
degradation firing transfers one token from $y$ to $\overline y$.  Therefore
the capacity identity is invariant.

The initial marking is
\[
\begin{array}{c|cccccc}
 &x_1&x_2&\txa_1&\txa_2&\cxr_1&\cxr_2\\ \hline
M_0&0&0&79&79&76&76
\end{array}
\]
with each complementary marking determined by its conservation equation.

\section{Translation rule: state change versus regulation}

An SPN is written
\[
  \mathcal N=(P,T,\Pre,\Post,M_0,a),
\]
where $P$ is the place set, $T$ the transition set, and
$a_t(M)$ the marking-dependent propensity of transition $t$.  Transition $t$
is structurally enabled when
\[
  M(p)\geq \Pre(p,t)\qquad\text{for every }p\in P,
\]
and firing it produces
\[
  M'=M+C(:,t),\qquad C=\Post-\Pre.
\]

Two different roles must be kept separate.

\paragraph{The variable that changes.}
For production of $y$, put $\overline y$ in the pre-set and $y$ in the
post-set.  For degradation, reverse these arcs.  These are the only arcs that
change the marking.

\paragraph{A regulator that is observed but not changed.}
The implementation represents regulatory dependence by a self-loop: the same
place occurs in both $\Pre$ and $\Post$.  Its incidence is zero, so firing does
not alter the regulator.  Such a self-loop also imposes the structural test
$M(p)\geq1$.

The sign of the interaction determines which place is used:
\begin{itemize}
  \item for an \emph{activator} $A\rightarrow y$, use the positive place $A$
        in the self-loop; the kinetic factor is an increasing Hill function;
  \item for an \emph{inhibitor} $R\dashv y$, use the complementary place
        $\overline R$ in the self-loop; the kinetic factor is a decreasing Hill
        function evaluated from the positive marking $M(R)$.
\end{itemize}
This convention makes the graph display the regulatory sign.  The propensity,
not the self-loop alone, retains the quantitative Hill response.  In
particular, a self-loop is a stricter condition than a purely functional read
arc because it requires at least one token in its input place.

\section{Transition propensities}

Write $m_y=M(y)$ and define
\begin{align*}
 F_{\cxr}(m)&=1-\frac{m^{p_{13}}}
  {m^{p_{13}}+(p_{22}p_{14})^{p_{13}}},\\
 F_{\txa}(m_1,m_2)&=
 \frac{[\tfrac12(m_1+m_2)]^{p_{11}}}
 {[\tfrac12(m_1+m_2)]^{p_{11}}+(p_{23}p_{12})^{p_{11}}},\\
 G_{\txa}(m)&=1-\frac{m^{p_3}}
  {m^{p_3}+(p_{21}p_4)^{p_3}},\\
 G_{\cxr}(m)&=1-\frac{m^{p_5}}
  {m^{p_5}+(p_{21}p_6)^{p_5}}.
\end{align*}

The unsplit Xist production rate inherited from Eq.~\eqref{eq:xist} is
\[
 A_i(M)=p_{21}F_{\cxr}(m_{\cxr_i})
                 F_{\txa}(m_{\txa_1},m_{\txa_2}).
\]
To show explicitly which member of the shared tXA pool contributes to an
event, the implementation replaces this single transition by two transitions:
\[
 a_{j,i}^{x,+}(M)=
 \begin{cases}
 \displaystyle
 \frac{m_{\txa_j}}{m_{\txa_1}+m_{\txa_2}}A_i(M),
 &m_{\txa_1}+m_{\txa_2}>0,\\[6pt]
 0,&m_{\txa_1}+m_{\txa_2}=0,
 \end{cases}
 \qquad j\in\{1,2\}.
\]
This split preserves the original kinetics exactly:
\[
 a_{1,i}^{x,+}(M)+a_{2,i}^{x,+}(M)=A_i(M)
 \quad\text{when }m_{\txa_1}+m_{\txa_2}>0.
\]
It is a bookkeeping decomposition, not a new regulatory assumption.  The
fraction selects the relative contribution of $\txa_j$; the Hill factor still
depends on the total trans-acting pool, as in M\"utzel et al.

The remaining propensities are direct translations of
Eqs.~\eqref{eq:xist}--\eqref{eq:cxr}:
\begin{align*}
 a_i^{x,-}(M)&=0.1733m_{x_i},\\
 a_i^{\txa,+}(M)&=p_{23}G_{\txa}(m_{x_i}),&
 a_i^{\txa,-}(M)&=m_{\txa_i},\\
 a_i^{\cxr,+}(M)&=p_{22}G_{\cxr}(m_{x_i}),&
 a_i^{\cxr,-}(M)&=m_{\cxr_i}.
\end{align*}

\section{Complete transition definition}

Table~\ref{tab:transitions} defines all 14 transitions.  A plus sign denotes
production and a minus sign degradation.  Entries listed on both sides are
self-loops and therefore have zero net incidence.

\renewcommand{\arraystretch}{1.18}
\begin{longtable}{@{}p{2.3cm}p{3.4cm}p{3.4cm}p{4.6cm}@{}}
\caption{Pre-set, post-set, and propensity of every transition.}
\label{tab:transitions}\\
\toprule
Transition & Pre-set & Post-set & Propensity \\
\midrule
\endfirsthead
\toprule
Transition & Pre-set & Post-set & Propensity \\
\midrule
\endhead
$t^{x,+}_{1,1}$ & $\overline{x}_1,\txa_1,\overline{\cxr}_1$
 & $x_1,\txa_1,\overline{\cxr}_1$ & $a^{x,+}_{1,1}(M)$ \\
$t^{x,+}_{2,1}$ & $\overline{x}_1,\txa_2,\overline{\cxr}_1$
 & $x_1,\txa_2,\overline{\cxr}_1$ & $a^{x,+}_{2,1}(M)$ \\
$t^{x,-}_{1}$ & $x_1$ & $\overline{x}_1$ & $0.1733m_{x_1}$ \\
$t^{\txa,+}_{1}$ & $\overline{x}_1,\overline{\txa}_1$
 & $\overline{x}_1,\txa_1$ & $p_{23}G_{\txa}(m_{x_1})$ \\
$t^{\txa,-}_{1}$ & $\txa_1$ & $\overline{\txa}_1$ & $m_{\txa_1}$ \\
$t^{\cxr,+}_{1}$ & $\overline{x}_1,\overline{\cxr}_1$
 & $\overline{x}_1,\cxr_1$ & $p_{22}G_{\cxr}(m_{x_1})$ \\
$t^{\cxr,-}_{1}$ & $\cxr_1$ & $\overline{\cxr}_1$ & $m_{\cxr_1}$ \\
\midrule
$t^{x,+}_{1,2}$ & $\overline{x}_2,\txa_1,\overline{\cxr}_2$
 & $x_2,\txa_1,\overline{\cxr}_2$ & $a^{x,+}_{1,2}(M)$ \\
$t^{x,+}_{2,2}$ & $\overline{x}_2,\txa_2,\overline{\cxr}_2$
 & $x_2,\txa_2,\overline{\cxr}_2$ & $a^{x,+}_{2,2}(M)$ \\
$t^{x,-}_{2}$ & $x_2$ & $\overline{x}_2$ & $0.1733m_{x_2}$ \\
$t^{\txa,+}_{2}$ & $\overline{x}_2,\overline{\txa}_2$
 & $\overline{x}_2,\txa_2$ & $p_{23}G_{\txa}(m_{x_2})$ \\
$t^{\txa,-}_{2}$ & $\txa_2$ & $\overline{\txa}_2$ & $m_{\txa_2}$ \\
$t^{\cxr,+}_{2}$ & $\overline{x}_2,\overline{\cxr}_2$
 & $\overline{x}_2,\cxr_2$ & $p_{22}G_{\cxr}(m_{x_2})$ \\
$t^{\cxr,-}_{2}$ & $\cxr_2$ & $\overline{\cxr}_2$ & $m_{\cxr_2}$ \\
\bottomrule
\end{longtable}

In the Python implementation, these four split transitions are named
\texttt{t\_prod\_1\_x1}, \texttt{t\_prod\_2\_x1},
\texttt{t\_prod\_1\_x2}, and \texttt{t\_prod\_2\_x2}.  The mathematical
subscripts above identify first the tXA source $j$ and then the target Xist
allele $i$.

\section{Pre, Post, and incidence columns}

For a unit-weight transition, the matrix entries follow immediately from the
table.  For example, under the local place order
\[
 (x_1,\overline{x}_1,\txa_1,\overline{\txa}_1,
  \txa_2,\overline{\txa}_2,\cxr_1,\overline{\cxr}_1),
\]
the transition $t^{x,+}_{1,1}$ has
\[
 \Pre(:,t^{x,+}_{1,1})=
 \begin{pmatrix}0\\1\\1\\0\\0\\0\\0\\1\end{pmatrix},
 \qquad
 \Post(:,t^{x,+}_{1,1})=
 \begin{pmatrix}1\\0\\1\\0\\0\\0\\0\\1\end{pmatrix},
 \qquad
 C(:,t^{x,+}_{1,1})=
 \begin{pmatrix}1\\-1\\0\\0\\0\\0\\0\\0\end{pmatrix}.
\]
The tXA and complementary cXR self-loops appear in both $\Pre$ and $\Post$
but cancel in $C$.  Replacing $\txa_1$ with $\txa_2$ gives the column for
$t^{x,+}_{2,1}$, with exactly the same incidence vector.  The two transitions
therefore produce the same state change but have different propensities.

More generally, every production/degradation pair for species $y$ has
incidence
\[
 C(:,t_y^+)=e_y-e_{\overline y},\qquad
 C(:,t_y^-)=e_{\overline y}-e_y,
\]
which proves $M(y)+M(\overline y)$ is conserved.

\section{Why a variable is treated as activator or inhibitor}

The classification is not chosen from the instantaneous sign of a simulated
trajectory.  It comes from the regulatory architecture and the monotonicity
of the ODE term in the M\"utzel model:
\begin{itemize}
  \item If increasing regulator $r$ increases production of $y$, then
        $\partial P_y/\partial r>0$ and $r$ is an activator.  Here
        $\partial f(\txa)/\partial\txa_i\geq0$, so each $\txa_i$ activates
        both $x_1$ and $x_2$.
  \item If increasing $r$ decreases production of $y$, then
        $\partial P_y/\partial r<0$ and $r$ is an inhibitor.  Here
        $F_{\cxr}$ decreases with $\cxr_i$, so $\cxr_i$ inhibits $x_i$.
        Similarly, $G_{\txa}$ and $G_{\cxr}$ decrease with $x_i$, so Xist
        inhibits production of its local tXA and cXR.
  \item Cis/trans scope is determined biologically: $\cxr_i$ regulates only
        $x_i$, whereas the shared tXA term contains $\txa_1+\txa_2$ and acts
        on both Xist alleles.
\end{itemize}

The product $F_{\cxr}F_{\txa}$ in Eq.~\eqref{eq:xist} is a synergistic
combination: Xist production is large only when trans activation is strong
and cis repression is weak.  The SPN keeps this product in its propensity;
ordinary arcs alone do not encode it.

\section{Stochastic semantics and relation to the ODE}

At marking $M$, the implemented rate is
\[
 \widetilde a_t(M)=
 \ind\!\left\{M\geq\Pre(:,t)\right\}a_t(M).
\]
If optional guards are used, an additional factor $g_t(M)\in\{0,1\}$ is
multiplied into this expression.  The current cXR--tXA model uses Hill
propensities and structural enabling; its threshold guards are commented out.

Let $a_0(M)=\sum_t\widetilde a_t(M)$.  Gillespie's direct method samples
\[
 \tau=-\frac{\log U_1}{a_0(M)},\qquad U_1\sim\mathcal U(0,1),
\]
and selects transition $t$ with probability
\[
 \Pr(t\text{ fires next}\mid M)=\frac{\widetilde a_t(M)}{a_0(M)}.
\]
The expected infinitesimal drift is
\[
 \frac{d}{dt}\mathbb E[M\mid M(t)=M]
   =\sum_{t\in T}C(:,t)\widetilde a_t(M).
\]
Away from capacity boundaries, the two split Xist-production transitions have
the same incidence and their rates sum to $A_i(M)$.  Their contribution to the
$x_i$ drift is therefore exactly the positive term of
Eq.~\eqref{eq:xist}; degradation contributes $-0.1733m_{x_i}$.  The same
argument recovers Eqs.~\eqref{eq:txa} and \eqref{eq:cxr}.  Thus the SPN is a
discrete stochastic realization of the ODE reaction channels, with additional
finite-capacity enabling imposed by the complementary places.

\section{Modelling assumptions and points of interpretation}

\begin{enumerate}
  \item Tokens represent discrete molecular units on the scales
        $p_{21}$, $p_{22}$, and $p_{23}$.
  \item Hill functions are evaluated on positive-place markings, using the
        scaled thresholds of the M\"utzel stochastic model.
  \item Regulatory self-loops do not change a regulator, but they do impose a
        one-token enabling condition.  If regulation should remain defined at
        zero without structural disabling, a separate functional/read-arc
        relation is more appropriate.
  \item The two Xist-production channels are mutually alternative events.
        Splitting their rate proportionally to tXA abundance preserves the
        total rate and avoids an AND requirement that both tXA places be
        nonempty.
  \item The SPN contains the production and degradation channels of the six
        core variables.  The additional silencing-intermediate reactions used
        in the full stochastic simulations of M\"utzel et al.\ are outside
        the current implementation and are therefore not included here.
\end{enumerate}

\section{Summary}

The translation follows a systematic rule: decompose every ODE into
nonnegative production and degradation terms, create one stochastic event per
term, encode the actual state change with complementary-place token transfer,
and retain nonlinear regulation in the marking-dependent propensity.
Activators use their positive places as regulatory inputs, whereas inhibitors
use their complementary places.  In the present 14-transition model, each
Xist-production term is split according to the relative amounts of
$\txa_1$ and $\txa_2$; the split identifies the source contribution while
preserving exactly the total Xist-production rate of the M\"utzel equations.

\begin{thebibliography}{9}
\bibitem{mutzel2019}
V.~M\"utzel, I.~Okamoto, I.~Dunkel, M.~Saitou, L.~Giorgetti, E.~Heard,
and E.~G.~Schulz,
``A symmetric toggle switch explains the onset of random X inactivation in
different mammals,'' \emph{Nature Structural \& Molecular Biology},
vol.~26, pp.~350--360, 2019.
\href{https://doi.org/10.1038/s41594-019-0214-1}
{doi:10.1038/s41594-019-0214-1}.

\bibitem{mutzel2019supp}
V.~M\"utzel et al., ``Supplementary Notes: Model Description,'' Supplementary
Information to \emph{A symmetric toggle switch explains the onset of random
X inactivation in different mammals}, 2019, Supplementary Note~2: The
cXR--tXA model.
\href{https://pure.mpg.de/rest/items/item_3158407_1/component/file_3158409/content}
{Supplementary information}.
\end{thebibliography}

\end{document}
